\clearpage
\renewcommand{\sectioncolor}{nano}
\section{The target: what exists, what is missing}\label{sec:target}
\markboth{The target domain}{}

Here I want to be realistic rather than inspiring. The pasted text is right that biology spans
extraordinary scales. What it does not do is say \emph{which mathematics already covers which band},
and \emph{where the coverage actually stops}.

\subsection{The scale ladder, with coverage}

\begin{center}
\begin{tikzpicture}[font=\footnotesize]
% --- spatial axis ---
\draw[slate!55,line width=1.2pt,-{Stealth[length=5pt]}] (0,0)--(16.4,0);
\node[anchor=east,font=\scriptsize\bfseries,text=slate,align=right] at (-0.15,0) {LENGTH};
\foreach \x/\l/\n in {1.0/{$10^{-10}$\,m}/{atoms}, 4.1/{$10^{-9}$\,m}/{molecules},
                      7.2/{$10^{-7}$\,m}/{complexes}, 10.3/{$10^{-6}$\,m}/{cells},
                      13.4/{$10^{-3}$\,m}/{tissues}, 15.9/{$1$\,m}/{organism}}{
  \draw[slate!55,line width=0.8pt] (\x,-0.12)--(\x,0.12);
  \node[font=\scriptsize,above=3pt] at (\x,0.1) {\l};
  \node[font=\scriptsize,text=slate,below=3pt] at (\x,-0.1) {\n};}
% --- seams, drawn first so labels sit on top ---
\foreach \x in {4.1,7.3,10.35}{
  \draw[crimson!55,line width=1pt,dashed] (\x,-1.20)--(\x,-5.05);}
% --- coverage bands ---
\def\bandy{-1.55}
\foreach \a/\b/\c/\t/\r in {%
   0.5/4.6/quantum/{quantum chemistry, DFT, HEOM}/0,
   3.6/8.0/nano/{molecular dynamics, force fields}/-0.78,
   6.6/11.2/bio/{coarse-grained \& Markov state models}/-1.56,
   9.5/14.2/sand/{reaction networks, stochastic kinetics}/-2.34,
   12.6/16.3/ember/{reaction--diffusion, continuum}/-3.12}{
  \fill[\c,opacity=0.20,rounded corners=3pt] (\a,{\bandy+\r-0.30}) rectangle (\b,{\bandy+\r+0.30});
  \draw[\c,line width=0.9pt,rounded corners=3pt] (\a,{\bandy+\r-0.30}) rectangle (\b,{\bandy+\r+0.30});
  \node[font=\scriptsize,text=\c!70!black] at ({(\a+\b)/2},{\bandy+\r}) {\t};}
\node[font=\scriptsize,text=crimson,align=center,anchor=north] at (7.3,-5.15)
  {\textbf{the seams} --- these are stitched by hybrid schemes (QM/MM, multiscale coupling),\\
   \emph{not} bridged by a unified formalism. This is gap \textbf{(1)}.};
\end{tikzpicture}
\end{center}
\vspace{4pt}
\begin{center}\begin{minipage}{0.93\textwidth}\footnotesize\color{slate}
\textbf{Figure 3.}\; Which mathematics covers which length scale. The overlaps are real; the
\emph{transitions} are where the formalism changes character entirely, and nothing carries across
them automatically.
\end{minipage}\end{center}

\subsection{What already exists and works --- by name and date}

The pasted text says ``some of the ingredients already exist.'' Here are the specific results that
did the actual work, so a student can start reading tomorrow.

\begin{center}\small
\begin{tabular}{@{}>{\raggedright\arraybackslash}p{4.5cm}>{\raggedright\arraybackslash}p{11.7cm}@{}}
\toprule
\rowcolor{nano!13}
\textbf{\color{nano}Problem} & \textbf{\color{nano}The mathematics that solved it}\\
\midrule
Network topology $\Rightarrow$ dynamics, \emph{without} knowing rate constants &
 \textbf{Deficiency-zero and deficiency-one theorems} --- Horn \& Jackson 1972; Feinberg, 1972
 onwards. A structural property of the reaction graph guarantees existence, uniqueness and stability
 of equilibria. This is the outstanding example of mathematics tailored to molecular systems.\\
Why DNA replication is more accurate than binding energies allow &
 \textbf{Kinetic proofreading} --- Hopfield 1974 (PNAS); Ninio 1975. Error correction bought with
 energy dissipation.\\
Free energies from irreversible pulling &
 \textbf{Jarzynski equality} 1997 (PRL); \textbf{Crooks fluctuation theorem} 1999 (PRE).\\
Cost of precision in a driven system &
 \textbf{Thermodynamic uncertainty relation} --- Barato \& Seifert 2015 (PRL).\\
Small copy numbers, discrete molecules &
 \textbf{Chemical master equation}; \textbf{Gillespie} stochastic simulation algorithm, 1976--77.\\
Protein conformational dynamics from simulation &
 \textbf{Markov state models}; \textbf{persistent homology} (Edelsbrunner, Letscher \& Zomorodian
 2002; Carlsson 2009) applied to folding trajectories and to cryo-EM heterogeneity.\\
Strongly coupled, non-Markovian environments &
 \textbf{HEOM} --- Tanimura \& Kubo 1989. Numerically exact; expensive.\\
Composing biochemical modules &
 \textbf{Applied category theory} --- Baez \& Pollard, \textit{A compositional framework for reaction
 networks}, 2017; decorated cospans.\\
\bottomrule
\end{tabular}
\end{center}

\subsection{Where the mathematics genuinely runs out}

\begin{center}
\begin{tikzpicture}[font=\footnotesize,
  g/.style={rounded corners=4pt,draw=crimson,line width=1.2pt,fill=crimson!6,inner sep=9pt,
            text width=4.78cm,align=left,minimum height=4.55cm}]
 \node[g] (g1) at (-5.75,0)
   {\textbf{\color{crimson}GAP 1 --- Scale bridging}\\[4pt]
    QM/MM hybrids are \emph{stitching}, not a unified formalism.\\[4pt]
    There is no general theory of \textbf{what survives coarse-graining} for open, driven systems.\\[4pt]
    {\scriptsize The renormalisation group is the nearest relative --- and it works beautifully at
    equilibrium critical points, which is not where cells live.}};
 \node[g] (g2) at (0,0)
   {\textbf{\color{crimson}GAP 2 --- A variational principle far from equilibrium}\\[4pt]
    At equilibrium, \emph{minimise free energy} organises everything.\\[4pt]
    \textbf{There is no accepted analogue for driven systems.}\\[4pt]
    {\scriptsize Maximum entropy production and related proposals remain contested. This is a genuine
    open problem, not a gap in anyone's education.}};
 \node[g] (g3) at (5.75,0)
   {\textbf{\color{crimson}GAP 3 --- Non-Markovian open systems}\\[4pt]
    No clean general theory.\\[4pt]
    HEOM is powerful but is essentially numerically exact brute force.\\[4pt]
    {\scriptsize Precisely the regime a chromophore in a warm protein actually occupies --- which is
    why quantum biology is mathematically hard as well as experimentally hard.}};
\end{tikzpicture}
\end{center}

\begin{critbox}[title={The bottleneck that is not mathematics --- read this before choosing a problem}]
For most of molecular and cellular biology \textbf{right now}, the binding constraint is
\textbf{measurement}, not mathematics. We cannot yet watch single molecules inside a living cell with
sufficient spatial \emph{and} temporal resolution simultaneously. Inventing mathematics for data you
cannot take is the wrong order of operations.
\begin{center}
\begin{tikzpicture}[font=\scriptsize]
 \node[rounded corners=3pt,fill=sand!14,draw=sand!80!black,line width=0.9pt,inner sep=7pt,
       text width=15.4cm,align=center]
  {Which loops straight back to \textbf{Move 4}. If the constraint is measurement, then the highest-value
   contribution available is often \emph{an instrument} --- and the instrument will generate the
   mathematics, as the STM and AFM did.};
\end{tikzpicture}
\end{center}
\end{critbox}
